% Module 1 starter: Find and fix
% Academic Writing Using LaTeX, Central State University SemiAcademy
% Licence: CC BY 4.0 (https://creativecommons.org/licenses/by/4.0/)
%
% This short report has three planted problems:
%   1. one is an error that stops the compile,
%   2. one makes text vanish from the PDF without any error,
%   3. one is a package loaded in the wrong place (Overleaf reports an
%      error, but may still show a PDF).
% Find all three, fix them, and recompile until the log shows no errors
% and every sentence appears in the PDF.
% The measurements are invented for practice.
% Lines that start with % (like these) are comments. They never print.

\documentclass{article}

\title{Oxide Thickness After Dry Oxidation}
\author{Your Name}
\date{\today}

\begin{document}

\maketitle

\usepackage{amsmath}

\begin{abstract}
Silicon dioxide grown by dry oxidation is often used as an insulating
layer, so its thickness is checked after growth. We measured the oxide
on three wafers after one hour at 1000 degrees Celsius, at five points
on each wafer. The mean thickness was 50 nm, and on each wafer the
readings stayed within 3% of that wafer's mean, inside the 5% limit we set.
The next test uses a longer oxidation time.
\end{abstract}

\section{Introduction}
An oxide layer that is too thin can leak current, and one that is too
thick wastes furnace time. This report checks one furnace setting.

\section{Methods}
We oxidized three wafers together and measured each one at five points
with an ellipsometer, recording the \textbf{thickness in nanometres.
The mean thickness of a wafer is the average of its five readings:
\[
\bar{t} = \frac{1}{5} \sum_{i=1}^{5} t_i
\]

\section{Results}
The three wafers gave mean thicknesses of 49, 50 and 51 nm.
\begin{itemize}
  \item No reading was outside the limit.
  \item On each wafer, the center reading was slightly thicker than the
        edge readings.
\end{itemize}

\section{Conclusion}
On the three sampled wafers, one hour at 1000 degrees Celsius gave about
50 nm of oxide, with limited variation across the measured points. More
wafers and runs are needed before drawing a conclusion about the process.

\end{document}
